%-*-latex-*-
Polygonia is a high-end avant-garde architecture company.
Some of their designs are based on convex polygons.
For instance here's a convex polygon

\begin{python}
from latextool_basic import *
p = Plot()
p0 = (0,0)
p1 = (3,-1)
p2 = (4,2)
p3 = (2,3)
p4 = (-1,1)
p += Line(points=[p0,p1,p2,p3,p4,p0])
print(p)
\end{python}

\lq\lq Convex" means if you join two points of the polygon with a line segment,
the line segment is inside the polygon.
For course a fancy mansion has to be divided into rooms: bedrooms, banquet rooms,
study rooms, libraries, swimming pool rooms, cinema rooms, etc.
Polygonia is starting a new line of designs where
the rooms are created by building interior walls (walls inside the building)
joining vertices of a polygonal
building.
Here's an example where there are two interior walls for the above building:

\begin{python}
from latextool_basic import *
p = Plot()
p0 = (0,0)
p1 = (3,-1)
p2 = (4,2)
p3 = (2,3)
p4 = (-1,1)
p += Line(points=[p0,p1,p2,p3,p4,p0])
p += Line(points=[p0,p3], linecolor='red')
p += Line(points=[p0,p2], linecolor='red')
print(p)
\end{python}

and here's another:

\begin{python}
from latextool_basic import *
p = Plot()
p0 = (0,0)
p1 = (3,-1)
p2 = (4,2)
p3 = (2,3)
p4 = (-1,1)
p += Line(points=[p0,p1,p2,p3,p4,p0])
p += Line(points=[p1,p4], linecolor='red')
p += Line(points=[p1,p3], linecolor='red')
print(p)
\end{python}

Here's another wall design for a more complicated building:

\begin{python}
from latextool_basic import *
p = Plot()
p0 = (0,0)
p1 = (3,-1)
p2 = (4.5,0.8)
p3 = (4,2)
p4 = (2,3)
p5 = (0.5,2.5)
p6 = (-1,1)
p += Line(points=[p0,p1,p2,p3,p4,p5,p6,p0])
p += Line(points=[p1,p4], linecolor='red')
p += Line(points=[p1,p3], linecolor='red')
p += Line(points=[p4,p6], linecolor='red')
p += Line(points=[p4,p0], linecolor='red')
print(p)
\end{python}

In other words a design of interior walls is a collection of non--intersecting
line segments in the polygon building
such that each line segment joins two vertices of the polygon
and
the resulting rooms must be triangles.
The cost of interior walls is the sum of the lengths of the line segments
in the architectural design (think Pythagorus theorem).
The goal is to find the minimal cost for the interior
walls when given the vertices of the building.

(a) For a polygonal build with $n$ vertices,
how many different collections of interior wall design are there?
(SOLUTION PROVIDED.)

(b)
Find a suitable recursion that helps solve this problem.
(You have to find a function to optimize. Make sure you state the function clearly.)

(c)
Find and implement in C\texttt{++} a bottom-up DP solution to the problem.
Here's a sample execution
\begin{Verbatim}[frame=single,commandchars=\~\!\@]
~userinput!5@
~userinput!0.0 0.0 3.1 -1.0 4.2 2.0 2.1 3.1 -1.0 1.1@
12.7
{{0,2}, {0,3}}
\end{Verbatim}
The first input is the number of vertices of the polygon.
This is followed by the $(x,y)$--coordinates of the points listed in
counter-clockwise direction.
In the above, the polygonal is made up of points
$p_0 = (0.0, 0.0)$,
$p_1 = (3.1, -1.0)$,
$p_2 = (4.2, 2.0)$,
$p_3 = (2.1, 3.1)$, and
$p_4 = (-1.0, 1.1)$.
The output is 12.7 (the minimal cost of among all possible wall designs).
This is followed by the optimal wall design.
For instance in the above the first line wall is
\verb!{0,2}!, meaning the line segment $p_0, p_2$
and the second line segment is \verb!{0,3}! meaning $p_0, p_3$.
If there are two possible optimal wall designs, you must list all of them,
one complete wall design on each line.

(d) What is the runtime and the space complexity of your algorithm?
Write in the documentation of your main.cpp:
\begin{Verbatim}[frame=single,fontsize=\small]
// T(n) = O(?)
// SPACE(n) = O(?)

... your program ...
\end{Verbatim}

Solution and hint on the page next ...

\newpage
\textsc{WARNING ... incoming solution and hint}

(a)
Let $N(n)$ be the number of ways to triangulate an $n$--gon. Then
\[
N(n) = N(3)N(n-1) + N(4)N(n-2) + N(5)N(n-3) + \cdots + N(n-1)N(3)
\]
where $N(3) = 1$.
Let $N'(n) = N(n + 2)$.
Then the above becomes
\[
N'(n - 2) = N'(1)N'(n-3) + N'(2)N'(n-4) + N'(3)N'(n-5) + \cdots + N'(n-3)N'(1)
\]
Therefore
$N'(n - 2)$ is $C_{n-2}$, the $(n-2)$--th Catalan number.
Hence
\[
  N(n)
  =
  N'(n-2)
  = C_{n - 2}
  = \frac{1}{n - 2} \binom{2(n - 2 - 1)}{n - 2}
  = \frac{1}{n - 2} \binom{2(n - 3)}{n - 2}
\]
Asymptotically,
\[
  N(n)
  = C_{n - 2}
  = O(4^{n-2}/\sqrt{n-2})
\]


(c)
Consider this polygon as an example:
\begin{python}
from latextool_basic import *
p = Plot()
p0 = (0,0)
p1 = (3,-1)
p2 = (4.5,0.8)
p3 = (4,2)
p4 = (2,3)
p5 = (0.5,2.5)
p6 = (-1,1)

p += Line(points=[p0,p1,p2,p3,p4,p5,p6,p0])
#p += Line(points=[p1,p4], linecolor='red')
#p += Line(points=[p1,p3], linecolor='red')
#p += Line(points=[p4,p6], linecolor='red')
#p += Line(points=[p4,p0], linecolor='red')

from latexcircuit import *

X = POINT(x=p0[0], y=p0[1], label='$p_0$', anchor='east'); p += str(X)
X = POINT(x=p1[0], y=p1[1], label='$p_1$', anchor='north'); p += str(X)
X = POINT(x=p2[0], y=p2[1], label='$p_2$', anchor='west'); p += str(X)
X = POINT(x=p3[0], y=p3[1], label='$p_3$', anchor='south west'); p += str(X)
X = POINT(x=p4[0], y=p4[1], label='$p_4$', anchor='south'); p += str(X)
X = POINT(x=p5[0], y=p5[1], label='$p_5$', anchor='south east'); p += str(X)
X = POINT(x=p6[0], y=p6[1], label='$p_6$', anchor='east'); p += str(X)

print(p)
\end{python}
Let $C(0, 6)$ to be the cost of the polygon.
If I divide the polygon this way:
\begin{python}
from latextool_basic import *
p = Plot()
p0 = (0,0)
p1 = (3,-1)
p2 = (4.5,0.8)
p3 = (4,2)
p4 = (2,3)
p5 = (0.5,2.5)
p6 = (-1,1)

p += Line(points=[p0,p1,p2,p3,p4,p5,p6,p0])
#p += Line(points=[p1,p4], linecolor='red')
p += Line(points=[p0,p3], linecolor='red')
p += Line(points=[p3,p6], linecolor='red')
#p += Line(points=[p4,p0], linecolor='red')

from latexcircuit import *

X = POINT(x=p0[0], y=p0[1], label='$p_0$', anchor='east'); p += str(X)
X = POINT(x=p1[0], y=p1[1], label='$p_1$', anchor='north'); p += str(X)
X = POINT(x=p2[0], y=p2[1], label='$p_2$', anchor='west'); p += str(X)
X = POINT(x=p3[0], y=p3[1], label='$p_3$', anchor='south west'); p += str(X)
X = POINT(x=p4[0], y=p4[1], label='$p_4$', anchor='south'); p += str(X)
X = POINT(x=p5[0], y=p5[1], label='$p_5$', anchor='south east'); p += str(X)
X = POINT(x=p6[0], y=p6[1], label='$p_6$', anchor='east'); p += str(X)

print(p)
\end{python}
I can consider the cost $C(0, 3)$ and $C(3, 6)$
where
$C(0, 3)$ is the cost for polygon $p_0 p_1 p_2 p_3$
and 
$C(3, 6)$ is the cost for polygon $p_3 p_4 p_5 p_6$.


Note that the cost $C(0,6)$ is the cost of the diagonals in the
$p_0...p_6$ polygon.
Likewise the cost $C(0,3)$ is the cost of the diagonals in
the $p_0 ... p_3$ polygon.
The cost of the line $p_0p_3$ is not included in $C(0,3)$.
Likewise the cost of the line $p_3 p_6$ is not in the cost $C(3,6)$.
Therefore cost of triangulating $p_0...p_6$ where $p_0p_3p_6$ is chosen to be in the
triangulation is
\[
  C(0,3) + C(3,6) + c(p_0,p_3) + c(p_3,p_6)
\]
where $c(p_i, p_j)$ is the cost of the wall from $p_i$ to $p_j$.
Of course we have to go through all the possible cases:
\[
  C(0, 6)
  =
  \min
  \{
  C(0,k) + C(k, 6) + c(p_0, p_k) + c(p_k, p_6)
  \mid
  k = 1, 2, 3, ..., 5
  \}
\]
In general
\[
  C(i,j)
  =
  \min
  \{
  C(i,k) + C(k, j) + c(p_i, p_k) + c(p_k, p_j)
  \mid
  i + 1 \leq k \leq j - 1
  \}
\]
It should be clear that
\[
  C(i, i + 1) = 0
\]

(e) The runtime is $O(n^3)$ and the space is $O(n^2)$.
